\documentclass[11pt]{article}
\usepackage[margin=2.3cm]{geometry}
\usepackage{amsmath,amssymb,amsthm}
\usepackage[hidelinks]{hyperref}
\newtheorem{theorem}{Theorem}
\newtheorem{lemma}{Lemma}
\theoremstyle{remark}\newtheorem*{remark}{Remark}
\newcommand{\supp}{\operatorname{supp}}
\title{Complete polynomials with sparse squares:\\ an improvement of Verdenius' exponent via overlapping towers}
\author{Roberto Vetere\thanks{Computations, certification and first draft produced with an AI assistant (Claude, Anthropic) under the author's direction; see \S6. Code and data: \url{https://github.com/RobertoVetere/sparse-squares}.}}
\date{Draft, November 2026 --- not yet peer reviewed}
\begin{document}\maketitle

\begin{abstract}
For a polynomial $f$ let $\|f\|$ be its number of nonzero coefficients, and let $\hat Q(k)=\min\{\|f^2\| : f\in\mathbb R[x],\ \deg f=k-1,\ \|f\|=k\}$ (complete polynomials).
Verdenius (1949) proved $\hat Q(k)\ll k^{\log 8/\log 13}=k^{0.81071\ldots}$ along a sequence of $k$, which appears to be the best known upper bound.
We show $\hat Q(k)\le \tfrac{21}{9}\,k^{\log 9/\log 17}$, $\log 9/\log 17=0.77552\ldots$, for the infinite sequence $k=(20\cdot 17^K-4)/16$.
The construction is a tower $\prod_j p(x^{m^j})$ with \emph{overlapping} base $m<\deg p+1$, a real seed $p$ of degree $20$ whose square has a three-block support, certified by interval arithmetic, and a short $\ell$-adic argument proving that the tower stays complete.
The result concerns real coefficients only; for rational/integer coefficients we obtain no improvement over $\log 6/\log 9$.
\end{abstract}

\section{Background}
R\'enyi and Erd\H{o}s asked how sparse $f^2$ can be when $f$ has many terms. Erd\H{o}s \cite{E} proved $Q(k)<k^{1-c}$; Verdenius \cite{V} obtained the exponent $\log 8/\log 13$ for real (complete) polynomials and $\log 6/\log 9$ over $\mathbb Q$. The lower bound $Q(k)\gg\log k$ is due to Schinzel and Zannier \cite{S,SZ}, improving Schinzel. Coppersmith--Davenport \cite{CD} and Abbott \cite{A} studied the smallest complete examples (degree $12$). We are not aware of any improvement of Verdenius' upper bound; cf.\ \cite{SZ,M,EP}.

\section{The construction}
Let $p(x)=\sum_{i=0}^n a_ix^i$ have all $a_i\neq 0$, let $S=\supp(p^2)\subset[0,2n]$, and let $m\le n+1$. Put
\[P_K(x)=\prod_{j=0}^{K-1}p\bigl(x^{m^j}\bigr),\qquad \deg P_K=D_K=n\,\frac{m^K-1}{m-1}.\]
Since $P_K^2=\prod_j p^2(x^{m^j})$, we have $\supp(P_K^2)\subseteq\{\sum_{j<K}s_jm^j: s_j\in S\}$.
With $m=n+1$ this is Verdenius' setting ($P_K$ automatically complete). We take $m<n+1$, which creates more collisions in the digit sums but makes completeness non-trivial.

\paragraph{The seed.} $n=20$, $m=17$, and $p$ palindromic ($a_{20-i}=a_i$, $a_0=1$) with
\[\supp(p^2)=S=[0,6]\cup[17,23]\cup[34,40]=A+17E,\qquad A=[0,6],\ E=\{0,1,2\}.\]
By palindromicity it suffices to impose $[x^j]p^2=0$ for $j=7,\dots,16$: ten quadratic equations in $a_1,\dots,a_{10}$. (In general, $p\equiv\sqrt{L}\bmod x^{17}$ with $L=\sum_{j\le 6}[x^j]p^2$, and the equations express palindromicity of the truncated square-root series.) Total-degree homotopy gives $259$ isolated complex solutions, $42$ of them real with all $a_i\neq0$. One of them (to 20 digits) is
\begin{center}\small
$a_0..a_{10}=1,\ 1.0048830402228113379,\ 1.2538227425985270778,\ 0.46211051619273127379,\ 1.0221701367519609109,$\\
$0.37735544346794049561,\ 0.43548160112730375387,\ -1.3831004818966026266,\ 0.14704166870239142598,$\\
$0.99944106767275121901,\ -1.0658751424342707659$, \quad $a_{20-i}=a_i$.
\end{center}

\begin{lemma}[certificate]\label{cert}
The system above has exactly one solution in the box of radius $10^{-60}$ around the 110-digit refinement of these values; it is real, all $a_i\ne0$, and $\supp(p^2)=S$ exactly.
\end{lemma}
\noindent\emph{Proof.} Krawczyk's test in interval arithmetic (110 digits, \texttt{mpmath.iv}): $K(X)\subset\operatorname{int}X$. The intervals of all $a_i$ and of $[x^j]p^2$, $j\in S$, exclude $0$; zeros off $S$ hold exactly by construction and palindromicity. Script: \texttt{certify/krawczyk.py}. \qed

\begin{lemma}[support]\label{supp}
$\|P_K^2\|\le 7\cdot 9^{K-1}\cdot 3=21\cdot 9^{K-1}$.
\end{lemma}
\begin{proof}Write $s_j=\alpha_j+17\varepsilon_j$, $\alpha_j\in[0,6]$, $\varepsilon_j\in\{0,1,2\}$. Then $\sum_j s_j17^j=\alpha_0+\sum_{j=1}^{K-1}(\alpha_j+\varepsilon_{j-1})17^j+\varepsilon_{K-1}17^K$ and $\alpha_j+\varepsilon_{j-1}\in[0,8]$. The number of such values is at most $7\cdot9^{K-1}\cdot3$.\end{proof}

\begin{lemma}[unique minimal digit sum]\label{uniq}
Let $m-1\le n<2m$. Every $N\in[0,D_K]$ has a representation $N=\sum_{j<K}i_jm^j$ with $0\le i_j\le n$, and the one minimising $\sum_j i_j$ is unique.
\end{lemma}
\begin{proof}Existence is clear since $n\ge m-1$. Let $r=(i_j)\neq r'=(i'_j)$ both be minimal and let $j$ be the least index where they differ. Then $i_j\equiv i'_j\pmod m$ and $|i_j-i'_j|<2m$, so w.l.o.g.\ $i_j=i'_j+m$. If some $l>j$ has $i_l<n$, take the least such $t$ and replace $i_j\mapsto i_j-m$, $i_l\mapsto i_l+1-m$ $(j<l<t)$, $i_t\mapsto i_t+1$. All new digits lie in $[0,n]$, the value is unchanged, and the digit sum drops by $(t-j)(m-1)>0$, which contradicts minimality of $r$. Otherwise $i_l=n$ for all $l>j$, and then $\sum_{l>j}i'_lm^l=m^{j+1}+n\sum_{j<l<K}m^l$, which exceeds the largest value representable by digits $\le n$ in positions $j+1,\dots,K-1$. Contradiction.\end{proof}

\begin{lemma}[completeness]\label{dense}
Let $F=\mathbb Q(a_0,\dots,a_n)\subset\mathbb R$. Choose a prime $\mathfrak p$ of $F$ with $v_{\mathfrak p}(a_i)=0$ for all $i$ (all but finitely many primes qualify), and let $\ell$ be the rational prime below $\mathfrak p$. Then for $p_\ell(x)=p(\ell x)$, every coefficient of $P_K=\prod_{j<K}p_\ell(x^{m^j})$ in degrees $0,\dots,D_K$ is nonzero.
\end{lemma}
\begin{proof}$[x^N]P_K=\sum_r\bigl(\prod_j a_{i_j}\bigr)\ell^{\sum_j i_j}$ over representations $r$ as in Lemma~\ref{uniq}. The term of $r$ has valuation $v_{\mathfrak p}(\ell)\sum_ji_j$. By Lemma~\ref{uniq} exactly one term has minimal valuation, so the sum is nonzero.\end{proof}

\begin{theorem}
For every $K\ge1$ there is $f\in\mathbb R[x]$ of degree $D_K=20(17^K-1)/16$ with all $D_K+1$ coefficients nonzero and $\|f^2\|\le 21\cdot9^{K-1}$. Hence $\hat Q(k)\le\frac{21}{9}k^{\log9/\log17}$ for $k=D_K+1$.
\end{theorem}
\begin{proof}Take $f=P_K$ from Lemma~\ref{dense}. Scaling $x\mapsto\ell x$ does not change $\supp(p^2)$, so Lemma~\ref{supp} applies. Since $k=D_K+1\ge17^K$, we get $9^K\le k^{\log 9/\log 17}$.\end{proof}

\section{Remarks and limitations}
\begin{itemize}
\item The same scheme with $n=14$, $m=12$, $S=[0,4]\cup[12,16]\cup[24,28]$ gives $\log7/\log12\approx0.7831$ (also certified). More generally, the three-block family $n=3r+2$, $m=n-r/2$ gives $\log(r+3)/\log m$, which is minimal at $r=6$. Annealing over all symmetric patterns with $n\le32$ and all bases found nothing better.
\item The gain over Verdenius is asymptotic: our explicit bound $21\cdot 9^{K-1}$ falls below $k^{\log 8/\log 13}$ only from $K=7$ ($k\approx5\cdot10^8$).
\item \textbf{Rational coefficients:} none of the isolated real solutions we computed (all square symmetric systems with $n\le16$ and exponent below $0.816$, and $n=19,20$) is rational up to rescaling of $x$, apart from Verdenius' own degree-8 seed. The bound over $\mathbb Q$ ($\log6/\log9$, relevant to Erd\H{o}s Problem \#485 as stated over $\mathbb Q$) is not improved.
\item Extending the bound from this sequence to all $k$ should follow from Erd\H{o}s' interpolation argument \cite{E}; we have not written this out.
\item Heuristically, towers of a single seed whose zeros come from a square polynomial system seem limited to exponents around $0.77$. Exponent $<1/2$ (the Epoch AI ``FrontierMath: Open Problems'' target) would need a different source of cancellation.
\end{itemize}

\section{Verification}
Everything is reproducible from the repository: exhaustive and annealed pattern searches; carry-automaton growth rates (cross-checked by brute force); an independent tower check that does not import the search code (digit-sum counts $21\cdot9^{K-1}$ up to $K=6$ by direct enumeration, and $P_2$ complete with $\|P_2^2\|=189$ computed directly); and a brute-force check of Lemma~\ref{uniq} for $K\le4$.

\section{Disclosure}
The search, the certificates and this draft were produced with an AI assistant (Claude, Anthropic). The author has not yet had the arguments checked by an independent expert; comments and corrections are very welcome.

\begin{thebibliography}{9}\small
\bibitem{A} J.~Abbott, Sparse squares of polynomials, Math.\ Comp.\ 71 (2002), 407--413.
\bibitem{CD} D.~Coppersmith, J.~Davenport, Polynomials whose powers are sparse, Acta Arith.\ 58 (1991), 79--87.
\bibitem{E} P.~Erd\H{o}s, On the number of terms of the square of a polynomial, Nieuw Arch.\ Wisk.\ (2) 23 (1949), 63--65.
\bibitem{EP} T.~F.~Bloom, Erd\H{o}s Problem \#485, \url{https://www.erdosproblems.com/485}.
\bibitem{M} V.~Mantova, Algebraic equations with lacunary polynomials and the Erd\H{o}s--R\'enyi conjecture, Riv.\ Mat.\ Univ.\ Parma 7 (2016), 239--246.
\bibitem{R} A.~R\'enyi, On the minimal number of terms of the square of a polynomial, Hungarica Acta Math.\ 1 (1947), 30--34.
\bibitem{S} A.~Schinzel, On the number of terms of a power of a polynomial, Acta Arith.\ 49 (1987), 55--70.
\bibitem{SZ} A.~Schinzel, U.~Zannier, On the number of terms of a power of a polynomial, Rend.\ Lincei Mat.\ Appl.\ 20 (2009), 95--98.
\bibitem{V} W.~Verdenius, On the number of terms of the square and the cube of polynomials, Indag.\ Math.\ 11 (1949), 459--465.
\end{thebibliography}
\end{document}
